\section{Setup and conventions}

Throughout, $(\GG,\EE)$ is a finite simple graph with vertex set $\VV=\VV(\GG)$, and for
each $v\in\VV$ we are given a separable tracial von Neumann algebra $(M_v,\varphi_v)$ with
faithful normal trace $\varphi_v=\tau_v$. We let $(M,\varphi)$ denote the associated graph
product von Neumann algebra in the sense of Caspers--Fima \cite{CaFi17}; $\varphi$ is the
canonical trace and $J$ the modular conjugation associated with $\varphi$ on $L^2M$. For a
subgraph $\GG'\subseteq\GG$ we write $\Gamma(\GG')\subseteq M$ for the graph product
sub-von Neumann algebra generated by $\{M_v: v\in\VV(\GG')\}$; this is the image of the
graph product on $\GG'$ and carries a $\varphi$-preserving conditional expectation
$E_{\Gamma(\GG')}\colon M\to\Gamma(\GG')$.

We freely use the framework of \emph{relative proper proximality} for tracial von Neumann
algebras developed in \cite{DKEP23}. For a boundary piece $\XX$ (an $M$-$M$, $JMJ$-$JMJ$
hereditary $C^*$-subalgebra of $\B(L^2M)$ which is the closure of
$\overline{MJMJ\,\mathcal{X}_0\,JMJM}$ for a suitable corner system) we write
$\KKK^{\infty,1}_{\XX}(M)$ for the associated $\sharp$-class, $\tilde\S_{\XX}(M)$ for the
small-at-infinity algebra of \cite[\S8]{DKEP23}, and we use the characterisation of
relative proper proximality via the non-existence of an $M$-central state on
$\tilde\S_{\XX}(M)$ that restricts to a normal state on $M$ \cite[Lemma 8.5]{DKEP23}.
We use the notation $\iota^\sharp\colon\B(L^2M)\to\B(L^2M)^{\sharp*}_J$, and for a boundary
piece $\XX$ we write $q_{\XX}$ for the central support projection of $(\XX)^{\sharp*}_J$ in
$\B(L^2M)^{\sharp*}_J$. All boundary pieces arising below are those associated with
subalgebras of the form $\Gamma(\GG')$; the corresponding projection onto
$L^2(\Gamma(\GG'))\subseteq L^2M$ is denoted $e_{\GG'}$, and $\XX_{\Gamma(\GG')}$,
$\KKK^{\infty,1}_{\Gamma(\GG')}(M)$ the associated boundary piece and class as in
\cite{DKEP23}.

For the graph product Hilbert space we use the standard reduced-word decomposition
\[
L^2M \;=\; \mathbb{C}\Omega \;\oplus\; \bigoplus_{\substack{w=v_1\cdots v_n\\ \text{reduced}}}
\HH_{w},\qquad
\HH_{w}=\overline{L^2M_{v_1}^{o}\otimes\cdots\otimes L^2M_{v_n}^{o}},
\]
where $L^2M_v^{o}=L^2M_v\ominus\mathbb{C}\Omega$ and a word $w=v_1\cdots v_n$ in the vertices
is \emph{reduced} if it has no syllable cancellation/commutation reduction \cite{CaFi17}.
For a subgraph $\GG'$, $L^2(\Gamma(\GG'))$ is the closed span of $\mathbb{C}\Omega$ and those
$\HH_w$ with all letters of $w$ in $\VV(\GG')$. Recall the basic orthogonality:
$\HH_w\perp L^2(\Gamma(\GG'))$ whenever $w$ contains a letter $v_j\notin\VV(\GG')$.

We restate the main theorem.

\begin{maintheorem}\label{thm:main}
Let $\GG$ be a finite graph with $|\VV(\GG)|\geq 3$, and let $(M_v,\varphi_v)$ be separable
tracial for each $v\in\VV$. Assume $\GG$ is irreducible, i.e.\ there do not exist subgraphs
$\GG_1,\GG_2$ with $\VV(\GG_1)\cup\VV(\GG_2)=\VV(\GG)$ and $(v,u)\in\EE(\GG)$ for all
$v\in\VV(\GG_1)$, $u\in\VV(\GG_2)$ (with $\GG_1,\GG_2$ proper). If each $M_v$ contains a
trace-zero unitary, then $M$ is properly proximal.
\end{maintheorem}

\section{Intersection of boundary pieces}

We keep the notation of the introduction. Fix two subgraphs $\GG_1,\GG_2\subseteq\GG$, let
$M_i=\Gamma(\GG_i)$, and let $N=M_1\cap M_2=\Gamma(\GG_1\cap\GG_2)$. Write
$e_i=e_{\GG_i}$ for the orthogonal projection of $L^2M$ onto $L^2M_i$, and $e=e_{\GG_1\cap\GG_2}$
for the projection onto $L^2N$. Let $q_i=q_{\XX_{M_i}}$ and $q=q_{\XX_N}$ be the
corresponding central support projections in $\B(L^2M)^{\sharp*}_J$, so that
\[
q_i=\bigvee_{u,v\in\mathcal{U}(M)}\iota^\sharp(uJvJe_iJv^*Ju^*),\qquad
q=\bigvee_{u,v\in\mathcal{U}(M)}\iota^\sharp(uJvJeJv^*Ju^*).
\]

\begin{lemma}\label{lem:keyincl}
Let $\KK$ be a normal representation of $\B(L^2M)^{\sharp*}_J$ with
$\B(L^2M)^{\sharp*}_J\subseteq\B(\KK)$. Then
\[
\iota^\sharp(e_1)\,q_2\,\KK
=\overline{\iota^\sharp\big(e_1\,MJMJ\,e_2\big)\KK}
\subseteq \overline{\iota^\sharp\big(MJMJ\,e\big)\KK}=q\KK .
\]
\end{lemma}

\begin{proof}
That $q_2\KK=\overline{\iota^\sharp(MJMJe_2)\KK}$ and $q\KK=\overline{\iota^\sharp(MJMJe)\KK}$
is immediate from the definitions of $q_2,q$ and the fact that $\iota^\sharp(uJvJ)$ runs over
unitaries normalising the relevant algebras; multiplying by $\iota^\sharp(e_1)$ gives the
first equality, since $e_1\cdot MJMJe_2$ spans (a dense subset of) the indicated operators.

We decompose, using $M=M_1\oplus(M\ominus M_1)$ as operators acting by left and right
multiplication (here $M\ominus M_1$ denotes the elements of $M$ whose image in $L^2M$ is
orthogonal to $L^2M_1$):
\begin{align}
e_1 MJMJ e_2
&= e_1\,(M\ominus M_1)\,J(M\ominus M_1)J\,e_2
 + e_1\,M_1\,J(M\ominus M_1)J\,e_2 \notag\\
&\quad + e_1\,(M\ominus M_1)\,JM_1J\,e_2
 + e_1\,M_1\,JM_1J\,e_2. \label{eq:dec}
\end{align}
Since $e_1$ commutes with $M_1$ and with $JM_1J$ (both preserve $L^2M_1$) and
$e_1 M_1=M_1 e_1$, $e_1 JM_1J=JM_1J e_1$, we may rewrite the last three terms of
\eqref{eq:dec} as
\[
M_1\,e_1\,J(M\ominus M_1)J\,e_2,\qquad
JM_1J\,e_1\,(M\ominus M_1)\,e_2,\qquad
JM_1J\,M_1\,e_1 e_2 .
\]
For the last of these, $e_1 e_2$ is the projection onto $L^2M_1\cap L^2M_2=L^2N$, i.e.
$e_1e_2=e_2e_1=e$ (this identity is part of the standard structure: the intersection of the
word-subspaces over $\VV(\GG_1)$ and over $\VV(\GG_2)$ is the word-subspace over
$\VV(\GG_1)\cap\VV(\GG_2)$). Hence the fourth term is $JM_1J M_1\,e$, whose
$\iota^\sharp$-image applied to $\KK$ lies in $\iota^\sharp(MJMJe)\KK=q\KK$. It remains to
handle the first three terms.

\medskip
\emph{First term.} We claim
\begin{equation}\label{eq:claim1}
(e_1-e)\,(M\ominus M_1)\,J(M\ominus M_1)J\,e_2=0 .
\end{equation}
Taking adjoints (and using $J(M\ominus M_1)J$, $(M\ominus M_1)$ are $*$-closed up to the
trace structure; more precisely the adjoint statement is
$e_2\,(M\ominus M_1)\,J(M\ominus M_1)J\,(e_1-e)=0$, which is equivalent to
\eqref{eq:claim1}), it suffices to prove
\[
e_2\,(M\ominus M_1)\,J(M\ominus M_1)J\,(e_1-e)=0 .
\]
Now $(e_1-e)L^2M=L^2M_1\ominus L^2N$ is spanned by the $\HH_w$ with $w$ reduced, all letters
in $\VV(\GG_1)$, and at least one letter in $\VV(\GG_1)\setminus\VV(\GG_1\cap\GG_2)
=\VV(\GG_1)\setminus\VV(\GG_2)$. Applying $J(M\ominus M_1)J$ followed by $(M\ominus M_1)$
acts on $L^2M$ by right- and left-multiplication by elements whose reduced expansions use
only letters \emph{not} in $\VV(\GG_1)$, hence in particular not killing the existing letter
in $\VV(\GG_1)\setminus\VV(\GG_2)$: every resulting reduced word still contains a letter
$v_j\in\VV(\GG_1)\setminus\VV(\GG_2)\subseteq\VV(\GG)\setminus\VV(\GG_2)$. Such a word
$\HH_w$ is orthogonal to $L^2M_2$, whence $e_2$ annihilates it. This proves
\eqref{eq:claim1}. Consequently
\[
\iota^\sharp\big(e_1(M\ominus M_1)J(M\ominus M_1)Je_2\big)\KK
=\iota^\sharp\big(e(M\ominus M_1)J(M\ominus M_1)Je_2\big)\KK
\subseteq\iota^\sharp(MJMJe)\KK=q\KK .
\]

\medskip
\emph{Second term.} We claim $(e_1-e)\,J(M\ominus M_1)J\,e_2=0$. Indeed
$(e_1-e)L^2M$ is spanned by $\HH_w$ ($w$ reduced, all letters in $\VV(\GG_1)$, some letter
in $\VV(\GG_1)\setminus\VV(\GG_2)$); applying $J(M\ominus M_1)J$ (right multiplication by
elements with reduced words avoiding $\VV(\GG_1)$) produces words still containing a letter
in $\VV(\GG_1)\setminus\VV(\GG_2)$, hence orthogonal to $L^2M_2$, so $e_2$ kills them. (The
argument is the adjoint form, identical to that for the first term.) Therefore
$e_1 J(M\ominus M_1)J e_2=e\,J(M\ominus M_1)J\,e_2$, and so
\[
\iota^\sharp\big(M_1 e_1 J(M\ominus M_1)J e_2\big)\KK
=\iota^\sharp\big(M_1 e\, J(M\ominus M_1)J e_2\big)\KK
\subseteq\iota^\sharp(MJMJe)\KK=q\KK,
\]
using that $M_1 e=e M_1 e\subseteq MJMJe$ after absorbing $e$ as the chosen base projection
and $M_1\subseteq M$.

\medskip
\emph{Third term.} By the same orthogonality, $(e_1-e)(M\ominus M_1)e_2=0$, since
left-multiplying $(e_1-e)L^2M$ by elements of $(M\ominus M_1)$ keeps a letter in
$\VV(\GG_1)\setminus\VV(\GG_2)$, producing vectors orthogonal to $L^2M_2$. Hence
$e_1(M\ominus M_1)e_2=e(M\ominus M_1)e_2$ and
\[
\iota^\sharp\big(JM_1J\,e_1(M\ominus M_1)e_2\big)\KK
=\iota^\sharp\big(JM_1J\,e(M\ominus M_1)e_2\big)\KK
\subseteq\iota^\sharp(MJMJe)\KK=q\KK .
\]

Combining the four terms via \eqref{eq:dec}, we obtain
$\iota^\sharp(e_1 MJMJ e_2)\KK\subseteq q\KK$, as claimed.
\end{proof}

\begin{corollary}\label{cor:comm}
$[q_1,q_2]=0$ and $q_1q_2=q$.
\end{corollary}

\begin{proof}
Realise $\B(L^2M)^{\sharp*}_J$ faithfully and normally on a Hilbert space $\KK$ as above.
Since $q_2$ commutes with $\iota^\sharp(M)$ and $\iota^\sharp(JMJ)$ and
$q_1=\bigvee_{u,v\in\mathcal U(M)}\iota^\sharp(uJvJ)\,\iota^\sharp(e_1)\,\iota^\sharp(Jv^*Ju^*)$,
to prove $[q_1,q_2]=0$ it suffices to show $[\iota^\sharp(e_1),q_2]=0$.

By Lemma~\ref{lem:keyincl}, $\iota^\sharp(e_1)q_2\KK\subseteq q\KK$, so $\iota^\sharp(e_1)q_2$
has range inside $q\KK$; as $q$ is a projection,
$\iota^\sharp(e_1)q_2=q\,\iota^\sharp(e_1)q_2$. Moreover $q\le q_2$ (because
$\XX_N\subseteq\XX_{M_2}$, as $L^2N\subseteq L^2M_2$ gives $e\le e_2$ and hence
$\XX_N\subseteq\XX_{M_2}$, so $q\le q_2$). Therefore
$q\,\iota^\sharp(e_1)q_2=q_2 q\,\iota^\sharp(e_1)q_2=q_2\,\iota^\sharp(e_1)q_2$, i.e.
\begin{equation}\label{eq:onesided}
\iota^\sharp(e_1)q_2=q_2\,\iota^\sharp(e_1)q_2 .
\end{equation}
Taking adjoints in \eqref{eq:onesided} (recall $\iota^\sharp(e_1)$ and $q_2$ are self-adjoint
projections) yields $q_2\,\iota^\sharp(e_1)=q_2\,\iota^\sharp(e_1)q_2$, and comparing with
\eqref{eq:onesided} gives $\iota^\sharp(e_1)q_2=q_2\,\iota^\sharp(e_1)$, i.e.
$[\iota^\sharp(e_1),q_2]=0$. Hence $[q_1,q_2]=0$.

For the second statement, from $[\iota^\sharp(e_1),q_2]=0$ and \eqref{eq:onesided},
$\iota^\sharp(e_1)q_2=q_2\iota^\sharp(e_1)q_2=q\iota^\sharp(e_1)q_2$; using $q\le q_2$ and
$[\iota^\sharp(e_1),q_2]=0$ again,
\[
\iota^\sharp(e_1)q_2
= q\,\iota^\sharp(e_1)\,q_2
= q\,\iota^\sharp(e_1)\,q ,
\]
the last step because $q\le q_2$ and $q\,\iota^\sharp(e_1)q_2 q=q\iota^\sharp(e_1)q$, while
$q\iota^\sharp(e_1)q_2(1-q)$ has range in $q\KK\cap(1-q)\KK=0$. Therefore
\begin{align*}
q_1 q_2
&=\Big(\bigvee_{u,v}\iota^\sharp(uJvJe_1Jv^*Ju^*)\Big)q_2
 =\bigvee_{u,v}\iota^\sharp(uJvJ)\,\iota^\sharp(e_1)q_2\,\iota^\sharp(Jv^*Ju^*)\\
&=\bigvee_{u,v}\iota^\sharp(uJvJ)\,q\,\iota^\sharp(e_1)\,q\,\iota^\sharp(Jv^*Ju^*)
 = q\Big(\bigvee_{u,v}\iota^\sharp(uJvJe_1Jv^*Ju^*)\Big)q
 = q\,q_1\,q ,
\end{align*}
where in the last line we used that $q$ commutes with $\iota^\sharp(uJvJ)$ for all
$u,v\in\mathcal U(M)$ (it is a central support projection of an $M$-$M$, $JMJ$-$JMJ$
boundary piece). Finally $q\le q_1$ (since $e\le e_1$, exactly as for $q\le q_2$), so
$q q_1 q=q$. Hence $q_1q_2=q$, and by symmetry / $[q_1,q_2]=0$ also $q_2q_1=q$.
\end{proof}

\begin{corollary}\label{cor:intersection}
$\KKK^{\infty,1}_{M_1}(M)\cap\KKK^{\infty,1}_{M_2}(M)=\KKK^{\infty,1}_{N}(M)$, and
$q_{\XX_N}=q_{\XX_{M_1}}\wedge q_{\XX_{M_2}}=q_{\XX_{M_1}}q_{\XX_{M_2}}$.
\end{corollary}

\begin{proof}
Since $\KKK^{\infty,1}_{M_i}(M)=(\iota^\sharp)^{-1}\big((\XX_{M_i})^{\sharp*}_J\big)$ and
$(\XX_{M_i})^{\sharp*}_J=q_i\,\B(L^2M)^{\sharp*}_J\,q_i$, while
$(\XX_N)^{\sharp*}_J=q\,\B(L^2M)^{\sharp*}_J\,q$, we have
\[
(\XX_{M_1})^{\sharp*}_J\cap(\XX_{M_2})^{\sharp*}_J
= q_1\B(L^2M)^{\sharp*}_J q_1\cap q_2\B(L^2M)^{\sharp*}_J q_2
= (q_1\wedge q_2)\,\B(L^2M)^{\sharp*}_J\,(q_1\wedge q_2),
\]
the last equality holding because $q_1,q_2$ commute (Corollary~\ref{cor:comm}), so the
intersection of the two corners is the corner at $q_1q_2=q_1\wedge q_2$. By
Corollary~\ref{cor:comm}, $q_1\wedge q_2=q_1q_2=q$, hence this corner equals
$(\XX_N)^{\sharp*}_J$. Pulling back through $\iota^\sharp$ gives the stated equality of
classes.
\end{proof}

\section{Stability of relative proper proximality under intersections}

\begin{lemma}\label{lem:intersection-pp}
Suppose each $(M_v,\varphi_v)$ is separable and tracial. Let $\{\GG_i\}_{i=1}^s$ be subgraphs
of $\GG$, and suppose $M$ is properly proximal relative to $\Gamma(\GG_i)$ for each $i$. Then
$M$ is properly proximal relative to $\Gamma\big(\bigcap_i\GG_i\big)$.
\end{lemma}

\begin{proof}
Write $\GG_0=\bigcap_i\GG_i$. Suppose $M$ is \emph{not} properly proximal relative to
$\Gamma(\GG_0)$. By \cite[Lemma 8.5]{DKEP23}, there is an $M$-central state $\psi$ on
$\tilde\S_{\XX_{\Gamma(\GG_0)}}(M)$ whose restriction to $M$ is normal. Let
$q_0:=q_{\XX_{\Gamma(\GG_0)}}$ and $q^{(i)}:=q_{\XX_{\Gamma(\GG_i)}}$ be the corresponding
central support projections. By Corollary~\ref{cor:intersection} applied iteratively
(the $q^{(i)}$ pairwise commute, and their product equals the central support of the
intersection boundary piece),
\begin{equation}\label{eq:lattice}
q_0=\bigwedge_{i=1}^s q^{(i)},\qquad\text{equivalently}\qquad
q_0^\perp=\bigvee_{i=1}^s \big(q^{(i)}\big)^\perp .
\end{equation}

\emph{Case 1: $\psi(q_0)\neq 0$.} Then
$\psi':=\dfrac{1}{\psi(q_0)}\,\psi\circ\mathrm{Ad}(q_0)$ is a state. Since $q_0$ commutes with
$\iota^\sharp(M)$ and $\iota^\sharp(JMJ)$ and $\psi$ is $M$-central, $\psi'$ is $M$-central
and $\psi'|_M$ is normal. Moreover, for each $i$,
\[
\mathrm{Ad}(q_0)\,\B(L^2M)^{\sharp*}_J\subseteq q_0\,\B(L^2M)^{\sharp*}_J\,q_0
=(\XX_{\Gamma(\GG_0)})^{\sharp*}_J\subseteq(\XX_{\Gamma(\GG_i)})^{\sharp*}_J
\subseteq\tilde\S_{\XX_{\Gamma(\GG_i)}}(M),
\]
where the middle inclusion is $q_0\le q^{(i)}$ from \eqref{eq:lattice}. Hence $\psi'$ defines
an $M$-central state on $\tilde\S_{\XX_{\Gamma(\GG_i)}}(M)$ that is normal on $M$,
contradicting (via \cite[Lemma 8.5]{DKEP23}) that $M$ is properly proximal relative to
$\Gamma(\GG_i)$.

\emph{Case 2: $\psi(q_0)=0$, i.e.\ $\psi(q_0^\perp)=1$.} By \eqref{eq:lattice},
$q_0^\perp=\bigvee_i (q^{(i)})^\perp$, so since $\psi(q_0^\perp)=1$ there exists $i$ with
$\psi\big((q^{(i)})^\perp\big)\neq 0$. Set $p:=(q^{(i)})^\perp$ and
$\psi'':=\dfrac{1}{\psi(p)}\,\psi\circ\mathrm{Ad}(p)$, a state. As above $p$ commutes with
$\iota^\sharp(M)$ and $\iota^\sharp(JMJ)$, so $\psi''$ is $M$-central with $\psi''|_M$ normal.
We claim $\psi''$ restricts to a state on $\tilde\S_{\XX_{\Gamma(\GG_i)}}(M)$. Indeed,
\[
\mathrm{Ad}(p)\big(\tilde\S_{\XX_{\Gamma(\GG_0)}}(M)\big)
\subseteq p\,\B(L^2M)^{\sharp*}_J\,p\,\cap\,\iota^\sharp(JMJ)'
\subseteq\tilde\S_{\XX_{\Gamma(\GG_i)}}(M),
\]
the final inclusion holding because $p=(q^{(i)})^\perp$ and elements of
$p\,\B(L^2M)^{\sharp*}_J\,p$ commuting with $\iota^\sharp(JMJ)$ which are, modulo
$(\XX_{\Gamma(\GG_i)})^{\sharp*}_J=q^{(i)}\B(L^2M)^{\sharp*}_J q^{(i)}$, supported on
$p=1-q^{(i)}$ lie in the small-at-infinity algebra
$\tilde\S_{\XX_{\Gamma(\GG_i)}}(M)$ by definition (which contains
$\big(q^{(i)}\big)^\perp\B(L^2M)^{\sharp*}_J\big(q^{(i)}\big)^\perp\cap\iota^\sharp(JMJ)'$
together with the boundary piece itself). Thus $\psi''$ is an $M$-central state on
$\tilde\S_{\XX_{\Gamma(\GG_i)}}(M)$, normal on $M$, again contradicting relative proper
proximality with respect to $\Gamma(\GG_i)$ via \cite[Lemma 8.5]{DKEP23}.

In both cases we reach a contradiction, so $M$ is properly proximal relative to
$\Gamma(\GG_0)$.
\end{proof}

\section{Relative proper proximality of amalgamated free products}

\begin{lemma}\label{lem:afp-pp}
Let $M_1,M_2$ be separable tracial von Neumann algebras with a common expected subalgebra
$B$ (i.e.\ $B\subseteq M_i$ with trace-preserving conditional expectations $E_B^{(i)}$).
Suppose there exist unitaries $u_1,u_2\in\mathcal U(M_1)$ and $u_3\in\mathcal U(M_2)$ with
\[
E_B(u_1)=E_B(u_2)=E_B(u_3)=E_B(u_1^*u_2)=0 .
\]
Then the amalgamated free product $M:=M_1\ast_B M_2$ is properly proximal relative to $B$.
\end{lemma}

\begin{proof}
We use the standard notation for amalgamated free products. For $i=1,2$, recall the
projection $\lambda_i(e_B)\in\B(L^2M)$ onto the closed span of $L^2B$ together with all
reduced words whose first (leftmost) syllable lies in $M_{i+1}^o$ (indices mod $2$, where
$M^o_j=M_j\ominus B$); equivalently, $\lambda_i(e_B)$ is the projection onto the vectors
that "begin in $L^2M_{i+1}^o$ or $L^2B$." By \cite[Prop.~4.16]{toyosawa2025relative} one has
$\lambda_i(e_B)\in\S_B(M)\subseteq\tilde\S_B(M)$, and moreover the structural identity
\begin{equation}\label{eq:lambda-sum}
\lambda_1(e_B)+\lambda_2(e_B)=1+e_B ,
\end{equation}
where $e_B\in\B(L^2M)$ is the Jones projection onto $L^2B$.

Suppose, for contradiction, that $M$ is \emph{not} properly proximal relative to $B$. By
\cite[Lemma 8.5]{DKEP23} there is an $M$-central state $\varphi$ on $\tilde\S_B(M)$ with
$\varphi|_M$ normal.

\medskip
\emph{Step 1: $\varphi(\iota^\sharp(e_B))=0$.}
Suppose not, so $\varphi(\iota^\sharp(e_B))\neq0$. Let
$p:=\bigvee_{u\in\mathcal U(M)}\iota^\sharp(u e_B u^*)$, the central support projection of the
boundary piece $\XX_B$; it commutes with $\iota^\sharp(M)$. Then
$\varphi\circ\mathrm{Ad}(p)$ is a nonzero $M$-central positive functional on
$\B(L^2M)^{\sharp*}_J$ that does not vanish on $\iota^\sharp(Me_BM)$ (because
$\varphi(\iota^\sharp(e_B))\neq0$ and $p\ge\iota^\sharp(e_B)$). Pulling back, this produces a
normal-on-$M$, $M$-central state on the $C^*$-algebra generated by $M$, $JMJ$ and $Me_BM$
which does not vanish on the compact-over-$B$ ideal $\overline{Me_BM}$. By the
Ozawa--Popa local criterion \cite[Lemma 3.3]{OP10b} (equivalently
\cite[Thm.~2]{BH18amenableabsorption}), the existence of such a functional forces
$M\preceq_M B$ (intertwining of $M$ into $B$ inside $M$).
But $M\not\preceq_M B$: taking $w:=u_1u_2\in\mathcal U(M_1)$ we have, by freeness with
amalgamation and $E_B(u_1)=E_B(u_2)=E_B(u_1^*u_2)=0$, that $w^n$ is, for each $n$, a reduced
word of length $2n$ over $M_1^o$ (its syllables alternate appropriately and never lie in
$B$), so for all $x,y\in M$,
\[
\|E_B(x\,w^n\,y)\|_2\xrightarrow[n\to\infty]{}0 .
\]
By Popa's intertwining theorem this means $M\not\preceq_M B$, a contradiction. Hence
$\varphi(\iota^\sharp(e_B))=0$.

\medskip
\emph{Step 2: deriving the contradiction.}
By \eqref{eq:lambda-sum} and Step 1,
\[
\varphi\big(\iota^\sharp(\lambda_1(e_B))\big)+\varphi\big(\iota^\sharp(\lambda_2(e_B))\big)
=\varphi\big(\iota^\sharp(1+e_B)\big)=1+0=1 .
\]
Now we use the two operator inequalities
\begin{equation}\label{eq:ineqs}
u_1^*\,\lambda_1(e_B)\,u_1+u_2^*\,\lambda_1(e_B)\,u_2\ \le\ \lambda_2(e_B),
\qquad
u_3^*\,\lambda_2(e_B)\,u_3\ \le\ \lambda_1(e_B).
\end{equation}
We justify the first; the second is symmetric. The projection $\lambda_1(e_B)$ projects onto
$L^2B\oplus\bigoplus_{w:\ \text{first syllable}\in M_2^o}\HH_w$. For $u_1\in\mathcal U(M_1)$
with $E_B(u_1)=0$, the operator $u_1^*\lambda_1(e_B)u_1$ is the projection onto
$u_1^*\big(\mathrm{ran}\,\lambda_1(e_B)\big)$; since $u_1^*$ has trivial $B$-component, left
multiplication by $u_1^*$ sends a vector "beginning in $M_2^o$ or $B$" to a vector beginning
in $M_1^o$ (the new leftmost syllable is $u_1^*$ or its reduction, which is a nonzero element
of $M_1^o$), i.e.\ into $\mathrm{ran}\,\lambda_2(e_B)$. The same holds for $u_2$ with
$E_B(u_2)=0$. Moreover, because $E_B(u_1^*u_2)=0$, the two ranges
$u_1^*\,\mathrm{ran}\,\lambda_1(e_B)$ and $u_2^*\,\mathrm{ran}\,\lambda_1(e_B)$ are mutually
orthogonal: a vector $\xi$ common (up to the projections) to both would yield
$u_1^*\eta=u_2^*\eta'$ with $\eta,\eta'\in\mathrm{ran}\,\lambda_1(e_B)$, hence
$u_1^*u_2$ would map $\mathrm{ran}\,\lambda_1(e_B)$ nontrivially into itself preserving the
leading-syllable structure, contradicting $E_B(u_1^*u_2)=0$ (which forces $u_1^*u_2$ to
strictly increase, hence change, the leftmost-syllable type). Therefore the two summand
projections in the left-hand side of the first inequality of \eqref{eq:ineqs} are mutually
orthogonal subprojections of $\lambda_2(e_B)$, giving the inequality. The second inequality
follows identically from $E_B(u_3)=0$, $u_3\in\mathcal U(M_2)$.

Since $\varphi$ is $M$-central and $u_1,u_2\in\mathcal U(M)$, applying $\varphi$ to the first
inequality of \eqref{eq:ineqs} and using $\varphi(\iota^\sharp(u_j^* x u_j))
=\varphi(\iota^\sharp(x))$ for $x\in\tilde\S_B(M)$ (centrality), we obtain
\[
2\,\varphi\big(\iota^\sharp(\lambda_1(e_B))\big)
\le\varphi\big(\iota^\sharp(\lambda_2(e_B))\big).
\]
Likewise the second inequality with $u_3\in\mathcal U(M)$ gives
\[
\varphi\big(\iota^\sharp(\lambda_2(e_B))\big)
\le\varphi\big(\iota^\sharp(\lambda_1(e_B))\big).
\]
Combining,
\[
2\,\varphi\big(\iota^\sharp(\lambda_1(e_B))\big)
\le\varphi\big(\iota^\sharp(\lambda_2(e_B))\big)
\le\varphi\big(\iota^\sharp(\lambda_1(e_B))\big),
\]
which (all quantities being nonnegative) forces
$\varphi(\iota^\sharp(\lambda_1(e_B)))=0$ and hence
$\varphi(\iota^\sharp(\lambda_2(e_B)))=0$. But then
$\varphi(\iota^\sharp(\lambda_1(e_B)+\lambda_2(e_B)))=0$, contradicting the value $1$
computed above.

This contradiction shows that no such $M$-central, normal-on-$M$ state exists, i.e.\ $M$ is
properly proximal relative to $B$.
\end{proof}

\section{Proof of the main theorem}

\begin{proof}[Proof of Theorem~\ref{thm:main}]
For a vertex $v\in\VV(\GG)$ recall the standard subgraphs:
$\mathrm{link}(v)=\{u\in\VV : (u,v)\in\EE(\GG)\}$, $\mathrm{star}(v)=\{v\}\cup\mathrm{link}(v)$,
and $\GG\setminus\{v\}$ the full subgraph on $\VV\setminus\{v\}$. The fundamental
decomposition of graph products as amalgamated free products gives
\begin{equation}\label{eq:afp-decomp}
M=\Gamma(\GG\setminus\{v\})\ \ast_{\Gamma(\mathrm{link}(v))}\ \Gamma(\mathrm{star}(v)).
\end{equation}
Set $M_1:=\Gamma(\GG\setminus\{v\})$, $B:=\Gamma(\mathrm{link}(v))$, and
$M_2:=\Gamma(\mathrm{star}(v))$; these are tracial von Neumann algebras with trace-preserving
conditional expectations onto $B$, and \eqref{eq:afp-decomp} is an amalgamation over $B$.

\medskip
\emph{Producing the required unitaries.}
We verify the hypotheses of Lemma~\ref{lem:afp-pp} for the decomposition
\eqref{eq:afp-decomp}, for each fixed $v$.

First, $M_2=\Gamma(\mathrm{star}(v))$ contains a unitary $u_3$ with $E_B(u_3)=0$: take
$u_3\in M_v$ a trace-zero unitary (which exists by hypothesis). Since $v\notin\mathrm{link}(v)$,
$M_v$ is free from $B=\Gamma(\mathrm{link}(v))$ amalgamated trivially over $\mathbb C$ in the
relevant graph-product sense, and $E_B(u_3)=\tau_v(u_3)\,1=0$.

Next we claim $M_1=\Gamma(\GG\setminus\{v\})$ contains unitaries $u_1,u_2$ with
$E_B(u_1)=E_B(u_2)=E_B(u_1^*u_2)=0$. Note $B=\Gamma(\mathrm{link}(v))$ is the graph product
over $\VV(\mathrm{link}(v))=\VV(\mathrm{star}(v))\setminus\{v\}$, a subgraph of
$\GG\setminus\{v\}$.

Suppose first that there exist two distinct vertices $a,b\in\VV(\GG\setminus\{v\})\setminus
\VV(\mathrm{link}(v))$ (i.e.\ $a,b\ne v$ are not adjacent to $v$). Pick trace-zero unitaries
$x\in M_a$, $y\in M_b$. If $a$ and $b$ are non-adjacent, set $u_1:=x$, $u_2:=y$; then
$E_B(u_1)=E_B(u_2)=0$ (each letter lies outside $\mathrm{link}(v)$ and is centered) and
$u_1^*u_2=x^*y$ is a reduced word of length $2$ in letters outside $\mathrm{link}(v)$, hence
$E_B(u_1^*u_2)=0$. If $a$ and $b$ are adjacent, set $u_1:=x$, $u_2:=xy$ — but to keep both
centered choose $u_1:=x$ and $u_2:=y$ as before when this already works; if instead only one
such vertex $a$ exists, see below. More robustly, when there are at least two vertices
$a,b\notin\mathrm{link}(v)\cup\{v\}$ we may always produce $u_1,u_2$: take a trace-zero
unitary $x\in M_a$ and a trace-zero unitary $z\in M_a$ with $\tau_a(x^*z)=0$ if
$\dim M_a$ is large enough, or use $a,b$: take $u_1:=x\in M_a$, $u_2:=y\in M_b$ trace-zero;
then $E_B(u_1)=E_B(u_2)=0$, and $u_1^*u_2=x^*y$ is reduced of length $2$ over
$\{a,b\}\subseteq\VV\setminus\VV(\mathrm{link}(v))$ regardless of adjacency of $a,b$ (a
length-two reduced word has its letters outside $\mathrm{link}(v)$ so its $B$-component
vanishes), hence $E_B(u_1^*u_2)=0$.

It remains to treat the case where there is exactly one vertex
$u\in\VV(\GG\setminus\{v\})$ with $u\notin\mathrm{link}(v)$ (such $u$ exists, for otherwise
every vertex other than $v$ is adjacent to $v$, contradicting irreducibility with
$\GG_1=\{v\}$, $\GG_2=\VV\setminus\{v\}$, since $|\VV|\ge 3$ makes both parts nonempty and the
join condition would hold). With exactly one such $u$, every vertex in
$\VV\setminus\{u,v\}$ is adjacent to $v$, i.e.\ lies in $\mathrm{link}(v)$.

\emph{Subcase (a): $u$ and $v$ are non-adjacent (which holds by definition of $u$) and there
is a vertex $w\in\VV\setminus\{u,v\}$ with $w\notin\mathrm{link}(u)$.} Then $M_u$ and $M_w$
are free (their generating vertices are non-adjacent), and $w\in\mathrm{link}(v)$. Take a
trace-zero unitary $x\in M_u$ and a trace-zero unitary $y\in M_w$, and set
$u_1:=x$, $u_2:=yx$. Since $u\notin\mathrm{link}(v)$ we have $E_B(u_1)=E_B(x)=0$; also
$u_2=yx$ is reduced of length $2$ (as $u,w$ non-adjacent), with the $u$-syllable $x$
centered and lying outside $\mathrm{link}(v)$, so the reduced word $yx$ contains a letter
($x\in M_u$) outside $B$, giving $E_B(u_2)=0$. Finally $u_1^*u_2=x^*yx$; here the leftmost
and rightmost syllables are $x^*,x\in M_u^o$ and the middle is $y\in M_w^o$; since $u,w$ are
non-adjacent this is a reduced word of length $3$ containing the letter $x\in M_u^o$ outside
$\mathrm{link}(v)$, whence $E_B(u_1^*u_2)=0$.

\emph{Subcase (b): every vertex $w\in\VV\setminus\{u,v\}$ is adjacent to $u$ (i.e.\
$\mathrm{link}(u)\supseteq\VV\setminus\{u,v\}$) and, as established, lies in
$\mathrm{link}(v)$.} We claim this contradicts irreducibility, unless we can still extract
unitaries; in fact in this subcase we still produce $u_1,u_2$ as follows. Consider the
vertex set partition $A=\{u,v\}$ and $C=\VV\setminus\{u,v\}$. Every $w\in C$ is adjacent to
both $u$ and $v$. Thus the only possibly-missing edge among "$A$ joined to $C$" is none: each
$a\in A$ is adjacent to each $c\in C$. Hence $\GG$ is the join of $A$ and $C$ along all edges
between them, i.e.\ $\VV(\GG_1)\cup\VV(\GG_2)=\VV$ with $\GG_1$ on $A$, $\GG_2$ on $C$ and
all cross edges present — contradicting irreducibility (note $C\ne\varnothing$ since
$|\VV|\ge 3$, and $A\ne\varnothing$). Therefore Subcase (b) cannot occur.

In all admissible cases we have produced $u_1,u_2\in\mathcal U(M_1)$ and
$u_3\in\mathcal U(M_2)$ with $E_B(u_1)=E_B(u_2)=E_B(u_3)=E_B(u_1^*u_2)=0$.

\medskip
\emph{Applying the lemmas.}
By Lemma~\ref{lem:afp-pp} applied to \eqref{eq:afp-decomp}, $M$ is properly proximal relative
to $B=\Gamma(\mathrm{link}(v))$, for every $v\in\VV(\GG)$.

By Lemma~\ref{lem:intersection-pp} with the family $\{\mathrm{link}(v)\}_{v\in\VV(\GG)}$, $M$
is properly proximal relative to $\Gamma\big(\bigcap_{v\in\VV(\GG)}\mathrm{link}(v)\big)$.

Finally, $\bigcap_{v\in\VV(\GG)}\mathrm{link}(v)=\varnothing$ as a vertex set: a vertex $w$
lying in $\mathrm{link}(v)$ for every $v\in\VV$ would in particular satisfy $w\in\mathrm{link}(w)$,
which is impossible in a simple graph (no loops). Hence
$\Gamma\big(\bigcap_v\mathrm{link}(v)\big)=\Gamma(\varnothing)=\mathbb C$, and being properly
proximal relative to $\mathbb C$ is precisely the statement that $M$ is properly proximal
\cite[\S8]{DKEP23}.

Therefore $M$ is properly proximal.
\end{proof}